\section{Introduction}
An object can remain between a robot's fingers while resting on a table; a
simulator can also generate a contact record across a positive gap with no
active force constraint. Rejecting both situations answers a literal
\emph{no-contact} question, not necessarily a \emph{no-external-load} question.
Conversely, observing external load does not by itself establish that the
object would be lost after a specified removal. These distinctions matter when
an evaluator is used to accept a candidate or judge a robot-program repair.
They do not imply that environmental contact is generally undesirable:
extrinsic manipulation deliberately exploits it~\citep{extrinsic2014}.

We study these distinctions as a measurement problem. The programs come from
earlier human--coding-agent heuristic iteration, but we do not test whether
agents outperform other developers. Previous frozen Allegro experiments found
37 versus 48 environment-clear 15-second successes among 360 scenes, with all
11 repair gains confined to spheres. Their literal contact exclusion motivates
re-examination, not an assumption that a less restrictive evaluator will
recover the excluded grasps. Historical counts and new experiments retain
separate data roles and model conditions.

\paragraph{Relation to existing evaluation.}
Temporal monitoring of manipulation is established: SafeManip evaluates
property-grounded temporal safety across tasks and environments
\citep{huang2026safemanip}. Grasp evaluators already guide selection and
refinement, including learned evaluators with real-world transfer
\citep{lum2025getagrip}. Task-oriented wrench-space methods evaluate what
forces a grasp can supply~\citep{chen2023task}, while dynamic tool-use work
assesses perturbation resistance and downstream behavior
\citep{grasp_to_act2026}. DexGraspBench provides a multi-hand physical
evaluation setting~\citep{chen2025bodex,dexgraspbench}. Thus, a temporal
predicate, a force threshold, an evaluator-driven selection rule, or a replay
branch is not in itself a new contribution. Broader behavior-quality evaluation
also goes beyond a binary success flag~\citep{eval_actions2026}.

Our narrower question is whether combining geometry with a well-specified
full-wrench monitor yields a measurable benefit beyond that simpler monitor,
and what construct changes remain when it does not. We use free-object
controls, independently coded raw-evidence reconstruction, source-specific
external tasks, and exact-state paired interventions. A frozen candidate-pool
experiment then tests decisions on a separate future endpoint. The primary
comparisons are against an information-matched wrench baseline, not merely
against a known contact-record proxy.

The contribution is an executable separation of three physical evaluation
targets, an empirical account of their disagreements and limits, and a
precommitted test of incremental decision value. Negative outcomes are part of
that account. No new grasp controller, general causal identification theorem,
or real-world safety guarantee is proposed.

\begin{figure}[t]
\centering
\includegraphics[width=\linewidth]{contract_examples.pdf}
\caption{Simulation renderings of three registered free-object controls at the
fixed branch time. The middle scene changes a record-level verdict without
external load; the supported scene carries real load. Renderings are explanatory,
not numerical evidence. All labels are computed from complete traces. The
object is green, actuated hand pads blue, and environment gray.}
\label{fig:examples}
\end{figure}
